% This file is part of the Open Source project 'proTirone'
% (c) 2025 Karsten Reincke (https://github.com/protirone/protico.ltx)
% It is distributed under the terms of the creative commons license
% CC-BY-4.0 (= https://creativecommons.org/licenses/by/4.0/)

\selectlanguage{ngerman}
\begin{frame}  
  \frametitle{LF09:15:L3-to-L3-with-Hopping:Activity Reminder}
  \begin{center}
    \includegraphics[height=7.5cm]{\imgLf/arp-switch-router-hopping-activity.png}
  \end{center}
\end{frame}

\begin{frame}  
  \frametitle{LF09:15:L3-to-L3-with-Hopping:Sequence Reminder}

  \begin{center}
    \includegraphics[height=7.5cm]{\imgLf/arp-switch-router-hopping-sequence-condensed.png}
  \end{center}
\end{frame}

\begin{frame}  
 \frametitle{LF09:15:L3-to-L3-with-Hopping:Package Reminder}

 \begin{center}
   \includegraphics[height=7.5cm]{\imgGl/arp-package-structure.png}
 \end{center}
\end{frame}

\begin{frame}  
  \frametitle{LF09:15:L3-to-L3-with-Hopping:Embedded Packages}

  \begin{center}
    \includegraphics[width=12cm]{\imgLf/hopping-l3-2-l3.png}
  \end{center}
\end{frame}

\begin{frame}  
  \frametitle{LF09:15:L7-to-L7-with-Hopping}

  \begin{center}
    
\includegraphics[height=6cm]{\imgGl/question-252536-oca.png}
  \end{center}
\end{frame}

\begin{frame}  
  \frametitle{LF09:15:L7-to-L7-with-Hopping:Request}

  \begin{center}
    \includegraphics[height=7.5cm]{\imgLf/ports-sockets-request-activity.png}
  \end{center}
\end{frame}

\begin{frame}  
  \frametitle{LF09:15:L7-to-L7-with-Hopping:Response}

  \begin{center}
    \includegraphics[height=7.5cm]{\imgLf/ports-sockets-response-activity.png}
  \end{center}
\end{frame}

\begin{frame}  
  \frametitle{LF09:15:L7-to-L7-with-Hopping:Embedded Packages}

  \begin{center}
    \includegraphics[height=7.5cm]{\imgGl/ports-services-layers.png}
  \end{center}
\end{frame}

\begin{frame}  
  \frametitle{LF109:15:Packages,Ports\&Services:Portsystematik}

    \begin{itemize}
    \item \tiny{Systemports := 0 - 1023 ('well-known ports', registriert bei der IANA)}.
      \begin{itemize}
        \item \tiny{\texttt{  7} ECHO-Service} 
        \item \tiny{\texttt{ 20} FTP (Datenübertragung: nur TCP)} 
        \item \tiny{\texttt{ 21} FTP (Verbindungsaufbau und Steuerung: TCP, UDP)}
        \item \tiny{\texttt{ 22} SSH (TCP, UDP)}
        \item \tiny{\texttt{ 23} Telnet (unverschlüsseltes Textprotokoll, z. B. für Fernwartung)}
        \item \tiny{\texttt{ 25} SMTP (Simple Mail Transfer Protocol)}
        \item \tiny{\texttt{ 67} Bootstrap Protocol (auch von DHCP genutzt: nur UDP)}
        \item \tiny{\texttt{ 80} HTTP (TCP, UDP)}
        \item \tiny{\texttt{110} POP3 (Post Office Protocol 3 = Mailclient zu Mailserver)}
        \item \tiny{\texttt{123} NTP (Network Time Protocol)}
        \item \tiny{\texttt{143} IMAP (Internet Message Access Protocol = Mailclient zu Mailserver)}
        \item \tiny{\texttt{319} PTP (Precision Time Protocol):Events  \texttt{320} PTP:General Messages (nur UDP)}
        \item \tiny{\texttt{443} HTTPS (Hypertext Transfer Protocol over SSL/TLS: TCP, UDP)}
      \end{itemize}
    \item \tiny{Registrierte Ports := 1024 - 49151 (per Konvention / 'Gebrauchsrecht')}
      \begin{itemize}
        \item \tiny{\texttt{1434} MSSQL}
        \item \tiny{\texttt{1883} MQTT}
        \item \tiny{\texttt{3306} Mysql}
      \item \tiny{\texttt{8080} HTTP über Poxyserver als Vermittler zwischen Client und realem Server}
      \end{itemize}
    \end{itemize}

\end{frame}

\begin{frame}  
  \frametitle{LF09:15:Packages,Ports\&Services::Package System}

  \begin{center}
    \includegraphics[height=7.5cm]{\imgGl/package-family-tree.png}
  \end{center}
\end{frame}

\begin{frame}  
  \frametitle{LF09:15:Packages,Ports\&Services:Echo:Client}

  \begin{center}
    \begin{minipage}{0.9\textwidth}
      \inputminted[fontsize=\footnotesize,linenos]{python}{\snpLf/ex-0915-echo-client.py}
    \end{minipage}
  \end{center}

\end{frame}

\begin{frame}  
  \frametitle{LF09:15:Packages,Ports\&Services:Echo:Server}

  \begin{center}
    \begin{minipage}{0.9\textwidth}
      \inputminted[fontsize=\footnotesize,linenos]{python}{\snpLf/ex-0915-echo-server.py}
    \end{minipage}
  \end{center}

\end{frame}

\begin{frame}  
  \frametitle{LF09:15:Packages,Ports\&Services::TCP-Verbindung}

  \begin{center}
    \includegraphics[height=7.5cm]{\imgLf/tcp-ip-pkg-handshake.png}
  \end{center}
\end{frame}


